\documentclass[11pt,a4paper]{article}
\usepackage[utf8]{inputenc}
\usepackage[T1]{fontenc}
\usepackage[ngerman]{babel}
\usepackage[margin=2.4cm]{geometry}
\usepackage{booktabs}
\usepackage{longtable}
\usepackage{xcolor}
\usepackage[colorlinks=true,urlcolor=blue!60!black,linkcolor=black]{hyperref}
\usepackage{parskip}
\usepackage{enumitem}
\setlist{nosep,leftmargin=*}
\emergencystretch=3em

\newcommand{\code}[1]{\texttt{\small #1}}

\title{AWI-ESM3-4-2-veg-HR, CMIP7 piControl\\[2mm]
\large Stand nach \code{cli118}: deine \code{aicc}-0.2-Befunde zu \code{cli117}}
\author{Jan Streffing, AWI}
\date{11. September 2026}

\begin{document}
% babel-ngerman macht " aktiv, was in Code-Beispielen zuschlaegt. Die Shorthands
% werden erst bei \begin{document} scharfgeschaltet, also muss das hierhin.
\shorthandoff{"}
\maketitle

\section{Kurzfassung}

Danke fuer den Lauf mit \code{cc-plugin-aicc} 0.2 auf \code{cli117}. Alle
Befunde, die bei uns lagen, sind behoben und in \code{cli118} nachgemessen,
dem Jahr 1851 mit 529 Dateien. Fuenf davon haben wir nach zwei Nachkorrekturen
noch einmal geschrieben, dazu Abschnitt~\ref{sec:nach}.

\begin{center}
\begin{tabular}{lrr}
\toprule
\textbf{Befund aus deinem Lauf} & \textbf{\code{cli117}} & \textbf{\code{cli118}} \\
\midrule
\code{lon} unter \code{valid\_min = 0}               & 299 & 0 \\
\code{time2}/\code{time4} ohne \code{climatology}     &   3 & 0 \\
\code{g132}: zwei Saetze Koordinaten                  &   2 Variablen & 0 \\
\code{hur}: \code{missing\_value}/\code{\_FillValue} unterschiedlich getypt & 3 & 0 \\
Daten als \code{double}, obwohl der Data Request \code{real} verlangt & 227 & 0 \\
\code{long\_name} \code{Latitude}/\code{Longitude}   &  41 & 0 \\
\code{long\_name} \code{ap(k+1/2)}/\code{b(k+1/2)}   &   9 & 0 \\
\code{g239} im Vokabular und DRS                      &  12 & 0 \\
\code{plev7h}                                         &   5 & deaktiviert \\
\code{TIME001}                                        &  42 & 42 \\
\bottomrule
\end{tabular}
\end{center}

In \code{cli118} bleiben damit 43 HIGH ueber alle drei Checker: cf 0,
\code{wcrp\_cmip7} 42, \code{aicc} 1. Die 42 sind \code{TIME001} und haengen an
deinem \code{\#81}, der eine \code{aicc}-Befund ist \code{vegtype}. Beides kommt
in Abschnitt~\ref{sec:rest}.

\section{Laenge in $[0, 360]$}

Die Tabelle verlangt fuer \code{longitude} und fuer \code{vertices\_longitude} in
\code{CMIP7\_grids.json} \code{valid\_min = 0} und \code{valid\_max = 360}. In
\code{cli117} lagen 299 Dateien in $[-180, 180]$, betroffen waren
\code{g122} auf \code{ncells} sowie alle Dateien auf \code{g130} und \code{g132}.
Daneben lagen 161 \code{g122}-Dateien schon im anderen Bereich. Dasselbe
Gitterlabel kam also mit zwei Koordinatensaetzen, und das duerfte ein guter Teil
der Konsistenzbefunde zwischen den Datasets gewesen sein.

Hilfskoordinaten und Eckpunkte werden jetzt in den Bereich gebracht.
Dimensionskoordinaten bleiben unberuehrt, dort wuerde das Umklappen die
Monotonie zerstoeren, und \code{g129} lag ohnehin schon in $[0, 360]$. Zellen,
die ueber den Nullmeridian reichen, haben ihre Eckpunkte danach an beiden Enden
des Bereichs, so wie CMOR sie schreibt. Der cf-Test, ob die Koordinate in ihrer
Zelle liegt, rechnet die Laenge dafuer zurueck, das haben wir auch in eurem
\code{compliance-checker}-Checkout nachgesehen.

\section{\code{g132}: ein Satz Koordinaten}

Dein Hinweis, dass \code{tauuo}/\code{tauvo} ein anderes Minimum haben als
\code{hfx}/\code{hfy}, lag an zwei Quellen. \code{tauuo} und \code{tauvo} kommen mit
den Elementschwerpunkten aus XIOS, \code{hfx} und \code{hfy} bekamen unsere, als
Mittel der drei Eckknoten auf der Kugel. Unser Schritt hat die XIOS-Werte bisher
stehen lassen, sobald eine Datei ueberhaupt \code{lat} mitbrachte.

Die beiden Saetze wichen auf 3{,}96 Mio.\ Zellen voneinander ab, auf 1272
polnahen Zellen um bis zu 57{,}6 Grad in der Laenge. XIOS mittelt dort die
Eckpunkt-Laengen in Grad, und das geht in Polnaehe schief. Auf dem Elementgitter
ersetzen wir jetzt immer. Alle neun \code{g132}-Dateien in \code{cli118} haben
bitgleiche \code{lat}, \code{lon} und Eckpunkte.

\section{Datentyp nach dem Data Request}

Alle Gleitkommavariablen im CMIP7 Data Request sind \code{real}, also einfach
genau. 227 der 534 Dateien in \code{cli117} waren \code{double}, meist weil eine
Einheitenumrechnung oder eine Mittelung \code{float32}-Eingaben hochgestuft hatte.
Die Daten werden jetzt vor dem Schreiben auf \code{float32} gebracht, und
\code{missing\_value} und \code{\_FillValue} folgen dem Typ der Daten.

Bei \code{hur} lag noch eine zweite Ursache darunter. Die Feuchteachsen
\code{hur100p2pct} und \code{hur101pct} haben in der Koordinatentabelle den
\code{out\_name} \code{hur}, und unser Schritt, der Koordinaten auf den Tabellentyp
bringt, hat deshalb die Variable selbst wieder als \code{double} geschrieben, mit
\code{float32}-Fuellwerten. Das war genau die Typmischung, die der cf-Checker
gemeldet hat und an der \code{check\_appendix\_a} abgestuerzt ist. Der Schritt
fasst jetzt nur noch Koordinaten an. In \code{cli118} ist keine Variable mehr
\code{double}, und der Absturz ist weg.

\section{\code{time2} und \code{time4} als Klimatologie}

\code{pfull} \code{tclm} und \code{tas} \code{tmaxavg}/\code{tminavg} tragen jetzt
\code{time:climatology = "climatology\_bnds"} und kein \code{bounds} mehr. Fuer
\code{time4} spannt \code{climatology\_bnds} den jeweiligen Monat. Fuer
\code{time2} laufen die Grenzen vom Monatsanfang im ersten bis zum Monatsende
im letzten Jahr, das in den Akkumulator eingegangen ist. Den Mittelungszeitraum
aus \code{DReq\#32} legen wir damit so aus: alle cmorisierten Jahre des Laufs.
Im Testjahr fallen beide Formen zusammen.

Eine Kleinigkeit dabei: \code{climatology\_bnds} steht numerisch in den Einheiten
der Zeitachse auf der Platte, ohne eigene Attribute. xarray raeumt
\code{units} und \code{calendar} nur von der Variable ab, auf die \code{bounds}
zeigt, nicht von der unter \code{climatology}. Sonst haette cf Appendix A
\code{calendar} auf einer Datenvariable gemeldet.

Den Tippfehler \code{Satistics} im \code{long\_name} von \code{time4} uebernehmen
wir aus der Tabelle, der steht schon in deinem \code{CMIP7\_DReq\_Content\#31}.

\section{\code{long\_name} nach Tabelle}

Dimensionskoordinaten tragen jetzt \code{Latitude}/\code{Longitude} wie in
\code{CMIP7\_coordinate.json}. Die Hilfskoordinaten auf unstrukturierten Gittern
bleiben kleingeschrieben, dort gilt \code{CMIP7\_grids.json}. Die Formelterme der
Grenzen heissen \code{ap(k+1/2)} und \code{b(k+1/2)}.

\section{Gitterregister in \code{aicc}}

\code{aicc} fuehrt ein festes Register von \code{g100} bis \code{g236}. Seitdem
sind 17 Labels ins CV gekommen, die Spezialgitter \code{g010} bis \code{g014} und
\code{g237} bis \code{g248}, darunter unser \code{g239} und das LR-Gitter
\code{g238}. Wir uebergeben deshalb ein eigenes Register per
\code{-O aicc:grid\_config:<json>}: euer Register plus 15 Labels, der Typ jeweils
aus dem \code{grid\_type}, den EMD fuer das Label fuehrt. \code{g013} und
\code{g014} fehlen darin, fuer Punkte und Linien kennt \code{aicc} keine Topologie.
Die 137 vorhandenen Eintraege stimmen alle mit EMD ueberein.

Falls dir das recht ist, schreiben wir das als Issue bei \code{cc-plugin-aicc}
auf. Ein Register, das aus dem CV gelesen wird, wuerde es gar nicht erst
veralten lassen.

Mit dem Register wurde auch etwas sichtbar, das vorher an der Gitterpruefung
haengen blieb: \code{msftm} stand als \code{(time, lev, basin, lat)} auf der
Platte, die Tabelle verlangt \code{(time, basin, lev, lat)}. Behoben, die
Hauptvariable folgt jetzt der Dimensionsfolge des Data Request, sobald alle
ihre Dimensionen darauf abbilden.

\section{Globale Mittel auf \code{g010}}

Die 29 Variablen mit horizontalem Mittel lagen bisher auf \code{g190}. Nach der
Klaerung im CV-Komitee ist \code{g190} ein Box-Modell-Gitter
(\code{Essential-Model-Documentation\#1382}), und fuer globale Mittel gilt das
Spezialgitter \code{g010}. Umgestellt. Die \code{Variable Registry}-Befunde auf
\code{so} und \code{thetao} aus deinem Lauf tauchen bei uns nicht auf, beide Terme
stehen in \code{branded\_variable/} von \code{CMIP7-CVs}. Wir vermuten dort einen
aelteren CV-Stand auf deiner Seite.

\section{Was uebrig bleibt}
\label{sec:rest}

\begin{center}
\begin{tabular}{rll}
\toprule
\textbf{Anzahl} & \textbf{Befund} & \textbf{wartet auf} \\
\midrule
42 HIGH & \code{TIME001}, die halbe Zelle & \code{cc-plugin-wcrp\#81} \\
 1 HIGH & \code{vegtype} auf \code{landCoverFrac} & \code{CMIP7\_DReq\_Content\#29} \\
492 MEDIUM & \glqq Registry has no value for key 'cell\_measures'\grqq{} & \code{WCRP-universe\#334} \\
\bottomrule
\end{tabular}
\end{center}

\textbf{\code{TIME001}.} Laut Beschreibung von \code{\#81} addiert \code{TIME001}
bei subdaily und dekadischen Daten keine halbe Zelle mehr und prueft den
midpoint gegen die deklarierten Grenzen. Das sind genau unsere 42. Sobald
\code{\#81} gemergt ist, pruefen wir \code{cli118} damit.

\textbf{Die 492 MEDIUM.} Die sind neu, seit wir auf \code{cmip7@2.3.0} gewechselt
sind, und sie liegen nicht an den Dateien. Die Universe-Datenbank 3.2.0, die
esgvoc aus der Registry laedt, fuehrt fuer \code{known\_branded\_variable} kein
\code{cell\_measures} und \code{units} statt \code{cf\_units}, obwohl das JSON im
Repository beides noch hat. Ein Update von esgvoc 4.1.1 auf 6.1.0 aendert daran
nichts. Das ist \code{WCRP-universe\#334} und \code{cc-plugin-wcrp\#79}, und dein
\code{\#81} erwaehnt genau diese Kompatibilitaet (\code{units} vs.\
\code{cf\_units}). Wir erwarten also, dass die 492 mit \code{\#81} ebenfalls
verschwinden.

\textbf{Zwei Stellen, die wir im cf-Checker selbst sehen.} Sie stehen nicht in
der Tabelle, weil sie nicht an uns haengen:

\begin{itemize}
\item \code{formula\_terms} auf \code{lev\_bnds}, in deinem Lauf acht HIGH aus
  Appendix A. \code{cf\_1\_7.py} verlangt das Attribut auf den Grenzen einer
  parametrischen Vertikalkoordinate, wie CF 1.7 \S7.1 es vorschreibt, der
  Appendix-A-Test verbietet es dort. Eine Datei kann nicht beides erfuellen,
  wir behalten es daher.
\item Der Test fuer \code{cell\_methods} von Klimatologien ist auf
  \code{\^{}time:} verankert. Seit \code{time2}/\code{time4} richtige
  Klimatologien sind, faellt deshalb \code{area: mean time: maximum within days
  time: mean over days} als MEDIUM durch, obwohl CF andere Achsen davor zulaesst.
  Drei Dateien.
\end{itemize}

Beides wuerden wir bei \code{compliance-checker} aufschreiben, ein bestehendes
Issue dazu haben wir nicht gefunden.

\section{Die FileNotFoundErrors}

Das war nicht diese Woche. Gemeint sind die fuenf fehlenden Reports aus
\code{cli116} vom 18.\,August, beschrieben in der vierten Antwort:
\code{cchecker.py} ist dort beim Import gestorben, bevor er eine Datei
angefasst hat. Seitdem wiederholt der QC-Schritt den Aufruf, wenn kein Report
entsteht. \code{cli117} hat 534 Reports auf 534 Dateien geliefert,
\code{cli118} 529 auf 529. In den 63 Logs von \code{cli118} steht kein einziger
\code{FileNotFoundError}, und die Wiederholung hat nie ausgeloest.

Ob es die Lustre-Probleme waren, koennen wir nicht sagen, der Fehler war nicht
reproduzierbar.

\section{Wie gemessen wurde}
\label{sec:nach}

\code{cli118} lief am 10./11.\,September mit allen Korrekturen ausser zweien:
\code{hur} war noch \code{double} (die zweite Ursache oben) und \code{msftm} in der
alten Dimensionsfolge. Beides ist behoben, und die fuenf betroffenen Dateien
sind am 11.\,September in denselben Lauf neu geschrieben. Alle Zahlen oben
beziehen sich auf diesen Stand.

Checker-Stand: \code{compliance-checker} 6.1.0, \code{cc-plugin-wcrp}
\code{develop} \code{ba2de16}, \code{cc-plugin-aicc} 0.2.0 mit unserem Register,
esgvoc 6.1.0 mit \code{cmip7@2.3.0} und \code{universe@3.2.0}. Die
\code{plev7h}-Variablen sind wie angekuendigt deaktiviert, daher 529 statt 534
Dateien.

Gute Erholung im Urlaub. Das hier liegt bereit, wenn du zurueck bist.

\end{document}
